\documentclass[10pt,aspectratio=169]{beamer}
\usetheme{Madrid}
\usecolortheme{default}
\usepackage[utf8]{inputenc}
\usepackage[T1]{fontenc}
\usepackage[spanish,mexico]{babel}
\usepackage{amsmath,amssymb,mathtools}
\usepackage{tikz}
\usepackage{booktabs,array}
\setbeamertemplate{navigation symbols}{}
\setbeamertemplate{theorems}[numbered]
\newtheorem{observacion}{Observación}
\newtheorem{pregunta}{Pregunta guía}
\title{Cálculo Integral}
\subtitle{Integrales trigonométricas elementales y sustitución trigonométrica}
\author{Carlos Ernesto Martínez}
\institute{Tecnológico Nacional de México}
\date{}
\begin{document}

\begin{frame}\titlepage\end{frame}


%=========================================================
\section{Integrales de las funciones trigonométricas elementales}
%=========================================================

\begin{frame}{Antes de la sustitución trigonométrica}
Antes de estudiar sustituciones como
\[
x=a\sin\theta,\qquad x=a\tan\theta,\qquad x=a\sec\theta,
\]
conviene separar un problema anterior e independiente:
\begin{block}{Pregunta}
¿Cómo se integran las seis funciones trigonométricas elementales y cuáles de sus primitivas pueden reconstruirse sin memorizar una tabla?
\end{block}
Esta revisión será importante más adelante porque, después de una sustitución trigonométrica, tendremos que integrar expresiones que contienen senos, cosenos, tangentes, secantes, etc.
\end{frame}

\begin{frame}{Punto de partida: integrar es leer una derivada al revés}
Si conocemos
\[
F'(x)=f(x),
\]
entonces podemos leer la igualdad en sentido inverso:
\[
\boxed{\int f(x)\,dx=F(x)+C.}
\]
Por eso, antes de memorizar integrales trigonométricas, conviene recordar las derivadas básicas.
\begin{block}{Idea didáctica}
Primero preguntamos: ``¿reconozco al integrando como una derivada conocida?'' Sólo cuando la respuesta es no necesitamos transformar la expresión.
\end{block}
\end{frame}

\begin{frame}{Derivadas trigonométricas que ya conocemos}
\small
\[
\boxed{
\begin{aligned}
(\sin x)'&=\cos x, & (\cos x)'&=-\sin x,\\[1mm]
(\tan x)'&=\sec^2x, & (\cot x)'&=-\csc^2x,\\[1mm]
(\sec x)'&=\sec x\tan x, & (\csc x)'&=-\csc x\cot x.
\end{aligned}}
\]
Observe que aquí aparecen directamente seis integrandos:
\[
\cos x,\quad \sin x,\quad \sec^2x,\quad \csc^2x,
\quad \sec x\tan x,\quad \csc x\cot x.
\]
Por tanto, esas seis integrales no requieren una técnica nueva.
\end{frame}

\begin{frame}{Seis primitivas inmediatas}
\small
Leyendo las derivadas anteriores en sentido inverso:
\[
\boxed{
\begin{aligned}
\int \cos x\,dx&=\sin x+C,\\
\int \sin x\,dx&=-\cos x+C,\\
\int \sec^2x\,dx&=\tan x+C,\\
\int \csc^2x\,dx&=-\cot x+C,\\
\int \sec x\tan x\,dx&=\sec x+C,\\
\int \csc x\cot x\,dx&=-\csc x+C.
\end{aligned}}
\]
\begin{alertblock}{Atención a los signos}
Los signos negativos en las primitivas de $\sin x$, $\csc^2x$ y $\csc x\cot x$ provienen de las derivadas de $\cos x$, $\cot x$ y $\csc x$.
\end{alertblock}
\end{frame}

\begin{frame}{¿Qué integrales faltan?}
La tabla de derivadas no nos entrega inmediatamente
\[
\boxed{
\int\tan x\,dx,\qquad
\int\cot x\,dx,\qquad
\int\sec x\,dx,\qquad
\int\csc x\,dx.}
\]
No conviene tratarlas como cuatro fórmulas aisladas.
\vspace{2mm}
\begin{block}{Objetivo}
Vamos a reconstruir las cuatro hasta convertirlas en la estructura
\[
\boxed{\int\frac{u'}{u}\,dx=\ln|u|+C.}
\]
Así veremos que, en realidad, las cuatro tienen una misma raíz conceptual: el logaritmo natural.
\end{block}
\end{frame}

%---------------- Tangente ----------------
\subsection{Integral de la tangente}
\begin{frame}{Integral de $\tan x$: primer paso}
Queremos calcular
\[
I=\int\tan x\,dx.
\]
La tangente puede escribirse como cociente:
\[
\tan x=\frac{\sin x}{\cos x}.
\]
Entonces
\[
I=\int\frac{\sin x}{\cos x}\,dx.
\]
\begin{pregunta}
¿Qué función aparece en el denominador y dónde está su derivada?
\end{pregunta}
El denominador es $\cos x$ y
\[
(\cos x)'=-\sin x.
\]
El numerador contiene precisamente $\sin x$; sólo falta controlar el signo.
\end{frame}

\begin{frame}{Integral de $\tan x$: sustitución}
Tomamos
\[
\boxed{u=\cos x}.
\]
Entonces
\[
du=-\sin x\,dx
\qquad\Longrightarrow\qquad
\sin x\,dx=-du.
\]
Sustituyendo paso a paso,
\begin{align*}
I
&=\int\frac{\sin x}{\cos x}\,dx\\
&=\int\frac{-du}{u}\\
&=-\int\frac{du}{u}\\
&=-\ln|u|+C.
\end{align*}
Regresando a $x$,
\[
\boxed{\int\tan x\,dx=-\ln|\cos x|+C.}
\]
\end{frame}

\begin{frame}{¿Por qué también aparece $\ln|\sec x|$?}
Como
\[
\sec x=\frac{1}{\cos x},
\]
entonces
\begin{align*}
\ln|\sec x|
&=\ln\left|\frac{1}{\cos x}\right|\\
&=\ln 1-\ln|\cos x|\\
&=-\ln|\cos x|.
\end{align*}
Por tanto,
\[
\boxed{
\int\tan x\,dx
=-\ln|\cos x|+C
=\ln|\sec x|+C.}
\]
No son dos métodos distintos: son dos escrituras equivalentes de la misma primitiva.
\end{frame}

%---------------- Cotangente ----------------
\subsection{Integral de la cotangente}
\begin{frame}{Integral de $\cot x$: reconocer la misma estructura}
Ahora
\[
I=\int\cot x\,dx.
\]
Usamos
\[
\cot x=\frac{\cos x}{\sin x},
\]
de modo que
\[
I=\int\frac{\cos x}{\sin x}\,dx.
\]
Esta vez la situación es todavía más directa:
\[
(\sin x)'=\cos x.
\]
Así, el numerador es exactamente la derivada del denominador.
\end{frame}

\begin{frame}{Integral de $\cot x$: desarrollo completo}
Tomamos
\[
\boxed{u=\sin x},\qquad du=\cos x\,dx.
\]
Entonces
\begin{align*}
I
&=\int\frac{\cos x}{\sin x}\,dx\\
&=\int\frac{du}{u}\\
&=\ln|u|+C\\
&=\boxed{\ln|\sin x|+C}.
\end{align*}
Por tanto,
\[
\boxed{\int\cot x\,dx=\ln|\sin x|+C.}
\]
\begin{block}{Comparación}
$\tan x$ produce un signo negativo porque $(\cos x)'=-\sin x$; $\cot x$ no lo produce porque $(\sin x)'=\cos x$.
\end{block}
\end{frame}

\begin{frame}{$\tan x$ y $\cot x$: una sola idea}
Podemos ver ambas integrales sin hacer la sustitución explícita:
\[
\tan x
=\frac{\sin x}{\cos x}
=-\frac{(\cos x)'}{\cos x},
\]
\[
\cot x
=\frac{\cos x}{\sin x}
=\frac{(\sin x)'}{\sin x}.
\]
Por tanto, ambas son manifestaciones de
\[
\boxed{\int\frac{f'(x)}{f(x)}\,dx=\ln|f(x)|+C.}
\]
Ésta es la estructura que debemos aprender a reconocer; las fórmulas particulares salen de ella.
\end{frame}

%---------------- Secante ----------------
\subsection{Integral de la secante}
\begin{frame}{La secante: ¿por qué es menos evidente?}
Consideremos
\[
I=\int\sec x\,dx.
\]
Si intentamos una sustitución inmediata, aparece un problema:
\[
(\sec x)'=\sec x\tan x,
\]
pero en la integral sólo tenemos $\sec x$.
\vspace{2mm}
Tampoco escribir $\sec x=1/\cos x$ produce directamente $f'/f$, porque falta un factor $\sin x$.
\begin{alertblock}{Necesitamos una manipulación adicional}
Vamos a modificar la apariencia del integrando sin cambiar su valor, multiplicando por una forma conveniente de $1$.
\end{alertblock}
\end{frame}

\begin{frame}{¿Cómo se descubre el ``truco'' de la secante?}
Queremos fabricar un denominador cuya derivada aparezca en el numerador.
Recordemos
\[
(\sec x)'=\sec x\tan x,
\qquad
(\tan x)'=\sec^2x.
\]
Si sumamos las funciones,
\[
\frac{d}{dx}(\sec x+\tan x)
=\sec x\tan x+\sec^2x.
\]
Factorizando el lado derecho,
\[
\sec x\tan x+\sec^2x
=\sec x(\tan x+\sec x).
\]
¡Ésta contiene exactamente el factor $\sec x$ que queremos integrar!
\end{frame}

\begin{frame}{Secante: fabricar $u'/u$}
Multiplicamos por
\[
1=\frac{\sec x+\tan x}{\sec x+\tan x}.
\]
Entonces
\begin{align*}
I
&=\int\sec x
\frac{\sec x+\tan x}{\sec x+\tan x}\,dx\\[1mm]
&=\int
\frac{\sec^2x+\sec x\tan x}{\sec x+\tan x}\,dx.
\end{align*}
Ahora compare:
\[
\underbrace{\sec x+\tan x}_{u}
\qquad\Longrightarrow\qquad
\underbrace{\sec x\tan x+\sec^2x}_{u'}.
\]
La manipulación no fue mágica: fue diseñada para producir $u'/u$.
\end{frame}

\begin{frame}{Secante: terminar la sustitución}
Tomamos
\[
\boxed{u=\sec x+\tan x}.
\]
Entonces
\[
du=(\sec x\tan x+\sec^2x)\,dx.
\]
Por tanto,
\begin{align*}
I
&=\int\frac{du}{u}\\
&=\ln|u|+C\\
&=\boxed{\ln|\sec x+\tan x|+C}.
\end{align*}
Así,
\[
\boxed{\int\sec x\,dx=\ln|\sec x+\tan x|+C.}
\]
\end{frame}

\begin{frame}{Comprobación de la integral de la secante}
Siempre podemos verificar una primitiva derivando:
\small
\begin{align*}
\frac{d}{dx}\ln|\sec x+\tan x|
&=\frac{\sec x\tan x+\sec^2x}{\sec x+\tan x}\\
&=\frac{\sec x(\tan x+\sec x)}{\sec x+\tan x}\\
&=\sec x.
\end{align*}
Por lo tanto, la fórmula es correcta.
\begin{block}{Lección}
La derivación no sólo sirve para comprobar el resultado; también explica retrospectivamente por qué elegimos $\sec x+\tan x$.
\end{block}
\end{frame}

%---------------- Cosecante ----------------
\subsection{Integral de la cosecante}
\begin{frame}{La cosecante: repetir la estrategia, no la fórmula}
Queremos calcular
\[
I=\int\csc x\,dx.
\]
Recordemos
\[
(\csc x)'=-\csc x\cot x,
\qquad
(\cot x)'=-\csc^2x.
\]
Entonces
\[
\frac{d}{dx}(\csc x+\cot x)
=-\csc x\cot x-\csc^2x.
\]
Factorizando,
\[
=-\csc x(\cot x+\csc x).
\]
Esto sugiere fabricar nuevamente una razón $u'/u$.
\end{frame}

\begin{frame}{Cosecante: multiplicar por una forma de $1$}
Multiplicamos por
\[
1=\frac{\csc x+\cot x}{\csc x+\cot x}.
\]
Así,
\begin{align*}
I
&=\int\csc x
\frac{\csc x+\cot x}{\csc x+\cot x}\,dx\\[1mm]
&=\int
\frac{\csc^2x+\csc x\cot x}{\csc x+\cot x}\,dx.
\end{align*}
Pero
\[
\frac{d}{dx}(\csc x+\cot x)
=-(\csc^2x+\csc x\cot x).
\]
El numerador es el negativo de la derivada del denominador.
\end{frame}

\begin{frame}{Cosecante: sustitución y resultado}
Tomamos
\[
\boxed{u=\csc x+\cot x}.
\]
Entonces
\[
du=-(\csc^2x+\csc x\cot x)\,dx,
\]
y por ello
\begin{align*}
I
&=-\int\frac{du}{u}\\
&=-\ln|u|+C\\
&=\boxed{-\ln|\csc x+\cot x|+C}.
\end{align*}
Por tanto,
\[
\boxed{\int\csc x\,dx=-\ln|\csc x+\cot x|+C.}
\]
\end{frame}

\begin{frame}{Una segunda forma para la integral de $\csc x$}
Usamos la identidad
\[
\csc^2x=1+\cot^2x.
\]
Entonces
\begin{align*}
(\csc x-\cot x)(\csc x+\cot x)
&=\csc^2x-\cot^2x\\
&=1.
\end{align*}
Por tanto,
\[
\csc x-\cot x=\frac{1}{\csc x+\cot x}.
\]
Tomando logaritmos,
\[
\ln|\csc x-\cot x|=-\ln|\csc x+\cot x|.
\]
Así,
\[
\boxed{\int\csc x\,dx=\ln|\csc x-\cot x|+C.}
\]
\end{frame}

\begin{frame}{Comprobación de la forma alternativa}
Derivemos
\[
F(x)=\ln|\csc x-\cot x|.
\]
Entonces
\small
\begin{align*}
F'(x)
&=\frac{-\csc x\cot x+\csc^2x}{\csc x-\cot x}\\
&=\frac{\csc x(\csc x-\cot x)}{\csc x-\cot x}\\
&=\csc x.
\end{align*}
Así confirmamos nuevamente que
\[
\int\csc x\,dx=\ln|\csc x-\cot x|+C.
\]
Las dos formas logarítmicas representan la misma familia de primitivas.
\end{frame}

%---------------- Síntesis ----------------
\subsection{Síntesis y conexión con el logaritmo}
\begin{frame}{Resumen: las seis integrales inmediatas}
\footnotesize
\[
\boxed{
\begin{aligned}
\int\cos x\,dx&=\sin x+C,
&\int\sin x\,dx&=-\cos x+C,\\[1mm]
\int\sec^2x\,dx&=\tan x+C,
&\int\csc^2x\,dx&=-\cot x+C,\\[1mm]
\int\sec x\tan x\,dx&=\sec x+C,
&\int\csc x\cot x\,dx&=-\csc x+C.
\end{aligned}}
\]
Estas seis salen directamente de la tabla de derivadas. No requieren una técnica adicional.
\end{frame}

\begin{frame}{Resumen: las cuatro que reconstruimos}
\small
\[
\boxed{
\begin{aligned}
\int\tan x\,dx
&=-\ln|\cos x|+C
=\ln|\sec x|+C,\\[1mm]
\int\cot x\,dx
&=\ln|\sin x|+C,\\[1mm]
\int\sec x\,dx
&=\ln|\sec x+\tan x|+C,\\[1mm]
\int\csc x\,dx
&=-\ln|\csc x+\cot x|+C\\
&=\ln|\csc x-\cot x|+C.
\end{aligned}}
\]
La meta no es memorizar las cuatro líneas, sino recordar cómo reconstruirlas.
\end{frame}

\begin{frame}{¿Por qué aparece siempre el logaritmo?}
Recordemos la estructura fundamental
\[
\boxed{\int\frac{du}{u}=\ln|u|+C.}
\]
Y, en forma compuesta,
\[
\boxed{\int\frac{f'(x)}{f(x)}\,dx=\ln|f(x)|+C.}
\]
En los cuatro casos hicimos exactamente eso:
\begin{itemize}
\item $\tan x$ y $\cot x$: la estructura $f'/f$ ya estaba casi visible.
\item $\sec x$ y $\csc x$: tuvimos que fabricarla algebraicamente.
\end{itemize}
Por eso los logaritmos no aparecen por casualidad.
\end{frame}

\begin{frame}{Conexión con la definición del logaritmo natural}
Para $x>0$, podemos definir
\[
\boxed{\ln x=\int_1^x\frac{dt}{t}.}
\]
Por ello, la integral
\[
\int\frac{du}{u}
\]
es precisamente la estructura elemental que genera el logaritmo.
\vspace{2mm}
Las primitivas de $\tan x$, $\cot x$, $\sec x$ y $\csc x$ no son, por tanto, resultados aislados: todas terminan reduciéndose a esa misma estructura.
\end{frame}

\begin{frame}{Mapa conceptual de esta primera parte}
\[
\boxed{
\begin{array}{c}
\text{función trigonométrica}\\[1mm]
\downarrow\\
\text{¿es una derivada conocida?}\\[1mm]
\swarrow\qquad\searrow\\[-1mm]
\text{sí}\qquad\qquad\text{no}\\
\downarrow\qquad\qquad\downarrow\\
\text{integrar directamente}\qquad
\text{buscar o fabricar }f'/f\\[1mm]
\qquad\qquad\downarrow\\
\qquad\qquad\ln|f(x)|+C
\end{array}}
\]
Este patrón será útil cuando las sustituciones trigonométricas transformen una integral algebraica en una integral trigonométrica.
\end{frame}

\begin{frame}{Transición: ahora sí, sustitución trigonométrica}
Ya contamos con dos herramientas distintas:
\begin{enumerate}
\item sabemos integrar las funciones trigonométricas elementales;
\item sabemos reconocer la estructura $f'/f$ que conduce al logaritmo.
\end{enumerate}
Ahora estudiaremos una técnica diferente: usar identidades trigonométricas para simplificar expresiones algebraicas como
\[
\sqrt{a^2-x^2},\qquad
\sqrt{a^2+x^2},\qquad
\sqrt{x^2-a^2}.
\]
La pregunta cambia de ``¿cómo integro esta función trigonométrica?'' a
\[
\boxed{\text{¿cómo uso trigonometría para simplificar una expresión algebraica?}}
\]
\end{frame}


\begin{frame}{Objetivo de esta parte}
Al terminar esta sección debemos poder responder, sin depender de una tabla, cuatro preguntas:
\begin{enumerate}
\item ¿Qué expresión algebraica está dificultando la integral?
\item ¿Por qué conviene usar seno, tangente o secante?
\item ¿Cómo la identidad pitagórica elimina la raíz?
\item ¿Cómo regresamos de $\theta$ a $x$ mediante un triángulo rectángulo?
\end{enumerate}
\vspace{2mm}
La idea no es memorizar tres recetas, sino comprender el mecanismo
\[
\boxed{\text{estructura algebraica}\to\text{identidad}\to\text{sustitución}\to\text{triángulo}.}
\]
\end{frame}

\section{Idea central}
\begin{frame}{¿Por qué introducir trigonometría?}
Hasta ahora, una sustitución típica puede ser
\[
u=x^2+1,\qquad u=3x-2.
\]
En sustitución trigonométrica hacemos algo diferente: introducimos deliberadamente un ángulo $\theta$.
\vspace{2mm}
\begin{pregunta}
¿Por qué introducir una función trigonométrica si la integral original ni siquiera contiene trigonometría?
\end{pregunta}
\pause
Porque las identidades pitagóricas pueden transformar \textbf{sumas y diferencias de cuadrados} en \textbf{cuadrados perfectos}.
\end{frame}

\begin{frame}{Las tres identidades que gobiernan la técnica}
\[
\boxed{1-\sin^2\theta=\cos^2\theta}
\]
\[
\boxed{1+\tan^2\theta=\sec^2\theta}
\]
\[
\boxed{\sec^2\theta-1=\tan^2\theta}
\]
No son tres fórmulas desconectadas. Son exactamente las identidades que necesitamos para las estructuras
\[
a^2-x^2,\qquad a^2+x^2,\qquad x^2-a^2.
\]
\end{frame}

\begin{frame}{¿Qué queremos conseguir?}
Supongamos que aparece
\[
\sqrt{a^2-x^2}.
\]
La raíz sería mucho más sencilla si el radicando pudiera escribirse como un cuadrado:
\[
a^2-x^2=a^2\cos^2\theta.
\]
Entonces
\[
\sqrt{a^2-x^2}=\sqrt{a^2\cos^2\theta}=a|\cos\theta|.
\]
Con una elección adecuada del intervalo de $\theta$,
\[
\sqrt{a^2-x^2}=a\cos\theta.
\]
\begin{block}{Propósito esencial}
Convertir una suma o diferencia de cuadrados en un cuadrado perfecto mediante una identidad trigonométrica.
\end{block}
\end{frame}

\section{Cómo descubrir las sustituciones}
\begin{frame}{Caso algebraico 1: $a^2-x^2$}
No empezamos memorizando $x=a\sin\theta$. La deducimos.
\begin{align*}
a^2-x^2=a^2\left(1-\frac{x^2}{a^2}\right)=a^2\left[1-\left(\frac{x}{a}\right)^2\right].
\end{align*}
Ahora comparamos con
\[
1-\sin^2\theta=\cos^2\theta.
\]
Para que ambas expresiones tengan la misma estructura hacemos
\[
\frac{x}{a}=\sin\theta.
\]
Por tanto,
\[
\boxed{x=a\sin\theta}.
\]
\end{frame}

\begin{frame}{¿Qué gana la sustitución $x=a\sin\theta$?}
Si $x=a\sin\theta$, entonces
\begin{align*}
a^2-x^2=a^2-a^2\sin^2\theta=a^2(1-\sin^2\theta)=a^2\cos^2\theta.
\end{align*}
Por ello,
\[
\sqrt{a^2-x^2}=a|\cos\theta|.
\]
Además, al derivar la sustitución,
\[
\boxed{dx=a\cos\theta\,d\theta}.
\]
Observe que la misma función $\cos\theta$ aparece tanto en la raíz transformada como en $dx$. Esa compatibilidad es parte de la eficacia de la sustitución.
\end{frame}

\begin{frame}{Caso algebraico 2: $a^2+x^2$}
Factorizamos $a^2$:
\begin{align*}
a^2+x^2=a^2\left(1+\frac{x^2}{a^2}\right)=a^2\left[1+\left(\frac{x}{a}\right)^2\right].
\end{align*}
La identidad con la misma estructura es
\[
1+\tan^2\theta=\sec^2\theta.
\]
Por eso identificamos
\[
\frac{x}{a}=\tan\theta,
\]
y obtenemos
\[
\boxed{x=a\tan\theta},\qquad \boxed{dx=a\sec^2\theta\,d\theta}.
\]
\end{frame}

\begin{frame}{¿Qué gana la sustitución $x=a\tan\theta$?}
\begin{align*}
a^2+x^2=a^2+a^2\tan^2\theta=a^2(1+\tan^2\theta)=a^2\sec^2\theta.
\end{align*}
Así,
\[
\sqrt{a^2+x^2}=a|\sec\theta|.
\]
Si elegimos
\[
-\frac{\pi}{2}<\theta<\frac{\pi}{2},
\]
entonces $\cos\theta>0$ y, por tanto, $\sec\theta>0$:
\[
\boxed{\sqrt{a^2+x^2}=a\sec\theta}.
\]
\end{frame}

\begin{frame}{Caso algebraico 3: $x^2-a^2$}
Ahora escribimos
\begin{align*}
x^2-a^2=a^2\left(\frac{x^2}{a^2}-1\right)=a^2\left[\left(\frac{x}{a}\right)^2-1\right].
\end{align*}
Buscamos una identidad del tipo
\[
(\text{algo})^2-1=(\text{algo})^2.
\]
La identidad adecuada es
\[
\sec^2\theta-1=\tan^2\theta.
\]
Por ello,
\[
\boxed{x=a\sec\theta},\qquad
\boxed{dx=a\sec\theta\tan\theta\,d\theta}.
\]
\end{frame}

\begin{frame}{Las tres sustituciones, ahora deducidas}
\centering
\renewcommand{\arraystretch}{1.55}
\begin{tabular}{>{\centering\arraybackslash}m{3cm}>{\centering\arraybackslash}m{4.4cm}>{\centering\arraybackslash}m{3.7cm}}
\toprule
Estructura & Identidad & Sustitución\\
\midrule
$a^2-x^2$ & $1-\sin^2\theta=\cos^2\theta$ & $x=a\sin\theta$\\
$a^2+x^2$ & $1+\tan^2\theta=\sec^2\theta$ & $x=a\tan\theta$\\
$x^2-a^2$ & $\sec^2\theta-1=\tan^2\theta$ & $x=a\sec\theta$\\
\bottomrule
\end{tabular}
\vspace{4mm}
\begin{alertblock}{Idea que sí conviene recordar}
No memorice la última columna sin la columna central. La identidad es la razón de la sustitución.
\end{alertblock}
\end{frame}

\section{Valor absoluto e intervalos}
\begin{frame}{Un detalle que no debemos esconder: $\sqrt{u^2}=|u|$}
En general,
\[
\boxed{\sqrt{u^2}=|u|},
\]
no simplemente $u$.
Por ejemplo,
\[
\sqrt{a^2\cos^2\theta}=|a\cos\theta|.
\]
Si $a>0$,
\[
|a\cos\theta|=a|\cos\theta|.
\]
Para poder escribir $a\cos\theta$ necesitamos elegir un intervalo donde $\cos\theta\ge0$.
\begin{block}{Elección habitual para $x=a\sin\theta$}
\[
-\frac{\pi}{2}\le\theta\le\frac{\pi}{2}.
\]
En este intervalo, $\cos\theta\ge0$.
\end{block}
\end{frame}

\begin{frame}{¿Por qué restringimos $\theta$?}
La restricción cumple dos funciones:
\begin{enumerate}
\item permite eliminar correctamente valores absolutos, por ejemplo
\[
\sqrt{\cos^2\theta}=\cos\theta;
\]
\item hace que la función trigonométrica usada sea uno a uno y podamos recuperar el ángulo mediante su inversa, por ejemplo
\[
\sin\theta=\frac{x}{a}
\quad\Longrightarrow\quad
\theta=\arcsin\left(\frac{x}{a}\right).
\]
\end{enumerate}
Por eso la elección del intervalo no es un tecnicismo aislado: conecta la sustitución con las ramas principales de las funciones trigonométricas inversas.
\end{frame}

\section{El papel del triángulo}
\begin{frame}{¿Por qué aparece un triángulo al final?}
Supongamos que hicimos
\[
x=a\sin\theta.
\]
Después de integrar obtenemos una expresión en $\theta$, pero necesitamos una respuesta en $x$.
De la sustitución,
\[
\sin\theta=\frac{x}{a}
=\frac{\text{opuesto}}{\text{hipotenusa}}.
\]
Eso nos permite construir un triángulo rectángulo. El lado faltante se obtiene por Pitágoras.
\begin{block}{Función del triángulo}
El triángulo es un dispositivo de \textbf{reconstrucción}: traduce $\sin\theta$, $\cos\theta$, $\tan\theta$, $\sec\theta$, etc., nuevamente a expresiones algebraicas en $x$.
\end{block}
\end{frame}

\begin{frame}{Triángulo asociado a $x=a\sin\theta$}
\begin{columns}[T]
\column{0.47\textwidth}
De
\[
\sin\theta=\frac{x}{a}
\]
tomamos
\[
\text{opuesto}=x,\qquad \text{hipotenusa}=a.
\]
Por Pitágoras,
\[
\text{adyacente}=\sqrt{a^2-x^2}.
\]
Entonces
\[
\cos\theta=\frac{\sqrt{a^2-x^2}}{a},\qquad
\tan\theta=\frac{x}{\sqrt{a^2-x^2}}.
\]
\column{0.50\textwidth}
\centering
\begin{tikzpicture}[scale=.9]
\coordinate (A) at (0,0);\coordinate (B) at (4,0);\coordinate (C) at (4,3);
\draw[thick] (A)--(B)--(C)--cycle;
\draw (3.65,0)--(3.65,.35)--(4,.35);
\node at (2,-.35) {$\sqrt{a^2-x^2}$};\node[right] at (4,1.5) {$x$};\node[above left] at (2.05,1.65) {$a$};
\draw (.75,0) arc[start angle=0,end angle=36.87,radius=.75];\node at (.95,.32) {$\theta$};
\end{tikzpicture}
\end{columns}
\end{frame}

\begin{frame}{Triángulo asociado a $x=a\tan\theta$}
\begin{columns}[T]
\column{0.47\textwidth}
De
\[
\tan\theta=\frac{x}{a}
\]
tomamos
\[
\text{opuesto}=x,\qquad \text{adyacente}=a.
\]
Por Pitágoras,
\[
\text{hipotenusa}=\sqrt{a^2+x^2}.
\]
Por tanto,
\[
\sec\theta=\frac{\sqrt{a^2+x^2}}{a},\qquad
\sin\theta=\frac{x}{\sqrt{a^2+x^2}}.
\]
\column{0.50\textwidth}
\centering
\begin{tikzpicture}[scale=.9]
\coordinate (A) at (0,0);\coordinate (B) at (4,0);\coordinate (C) at (4,3);
\draw[thick] (A)--(B)--(C)--cycle;
\draw (3.65,0)--(3.65,.35)--(4,.35);
\node at (2,-.35) {$a$};\node[right] at (4,1.5) {$x$};\node[above left] at (2.1,1.65) {$\sqrt{a^2+x^2}$};
\draw (.75,0) arc[start angle=0,end angle=36.87,radius=.75];\node at (.95,.32) {$\theta$};
\end{tikzpicture}
\end{columns}
\end{frame}

\begin{frame}{Triángulo asociado a $x=a\sec\theta$}
\begin{columns}[T]
\column{0.47\textwidth}
De
\[
\sec\theta=\frac{x}{a}
=\frac{\text{hipotenusa}}{\text{adyacente}}
\]
tomamos
\[
\text{hipotenusa}=x,\qquad \text{adyacente}=a.
\]
Por Pitágoras,
\[
\text{opuesto}=\sqrt{x^2-a^2}.
\]
Así,
\[
\tan\theta=\frac{\sqrt{x^2-a^2}}{a},\quad
\sin\theta=\frac{\sqrt{x^2-a^2}}{x}.
\]
\column{0.50\textwidth}
\centering
\begin{tikzpicture}[scale=.9]
\coordinate (A) at (0,0);\coordinate (B) at (4,0);\coordinate (C) at (4,3);
\draw[thick] (A)--(B)--(C)--cycle;
\draw (3.65,0)--(3.65,.35)--(4,.35);
\node at (2,-.35) {$a$};\node[right] at (4,1.5) {$\sqrt{x^2-a^2}$};\node[above left] at (2.15,1.65) {$x$};
\draw (.75,0) arc[start angle=0,end angle=36.87,radius=.75];\node at (.95,.32) {$\theta$};
\end{tikzpicture}
\end{columns}
\end{frame}

\begin{frame}{Algoritmo razonado}
\begin{enumerate}
\item Identifique la expresión problemática: $a^2-x^2$, $a^2+x^2$ o $x^2-a^2$.
\item Factorice $a^2$ mentalmente para reconocer la identidad pitagórica.
\item Deduzca la sustitución trigonométrica apropiada.
\item Calcule $dx$.
\item Sustituya \textbf{toda} la integral, no sólo la raíz.
\item Simplifique usando la identidad pitagórica.
\item Integre respecto de $\theta$.
\item Construya el triángulo correspondiente.
\item Reescriba todas las funciones de $\theta$ en términos de $x$.
\item Simplifique y añada $C$.
\end{enumerate}
\end{frame}

\section{Ejemplo 1: $a^2-x^2$}
\begin{frame}{Ejemplo 1: planteamiento}
Calcular
\[
\boxed{I=\int\sqrt{a^2-x^2}\,dx},\qquad a>0.
\]
\begin{pregunta}
¿Qué estructura aparece dentro de la raíz?
\end{pregunta}
\[
a^2-x^2.
\]
La identidad correspondiente es
\[
1-\sin^2\theta=\cos^2\theta,
\]
por lo que elegimos
\[
\boxed{x=a\sin\theta}.
\]
\end{frame}

\begin{frame}{Ejemplo 1: transformar $dx$ y la raíz}
De
\[
x=a\sin\theta
\]
obtenemos
\[
\boxed{dx=a\cos\theta\,d\theta}.
\]
Por otra parte,
\begin{align*}
\sqrt{a^2-x^2}
&=\sqrt{a^2-a^2\sin^2\theta}\\
&=a\sqrt{1-\sin^2\theta}\\
&=a\sqrt{\cos^2\theta}\\
&=a|\cos\theta|.
\end{align*}
Elegimos $-\pi/2\le\theta\le\pi/2$, de modo que $\cos\theta\ge0$:
\[
\boxed{\sqrt{a^2-x^2}=a\cos\theta}.
\]
\end{frame}

\begin{frame}{Ejemplo 1: la integral ya está en $\theta$}
Sustituimos ambas piezas:
\begin{align*}
I
&=\int \underbrace{(a\cos\theta)}_{\sqrt{a^2-x^2}}
\underbrace{(a\cos\theta\,d\theta)}_{dx}\\
&=a^2\int\cos^2\theta\,d\theta.
\end{align*}
Ahora aparece una integral trigonométrica conocida. Utilizamos
\[
\cos^2\theta=\frac{1+\cos(2\theta)}{2}.
\]
Por tanto,
\[
I=\frac{a^2}{2}\int\bigl(1+\cos(2\theta)\bigr)\,d\theta.
\]
\end{frame}

\begin{frame}{Ejemplo 1: integrar con cuidado}
\begin{align*}
I
&=\frac{a^2}{2}\left(\theta+\frac12\sin(2\theta)\right)+C.
\end{align*}
¿Por qué aparece $1/2$?
\[
\int\cos(2\theta)\,d\theta=\frac12\sin(2\theta).
\]
Usamos ahora
\[
\sin(2\theta)=2\sin\theta\cos\theta
\]
para obtener
\[
\boxed{I=\frac{a^2}{2}\theta+
\frac{a^2}{2}\sin\theta\cos\theta+C.}
\]
Todavía no hemos terminado: la respuesta sigue escrita en $\theta$.
\end{frame}

\begin{frame}{Ejemplo 1: reconstrucción mediante el triángulo}
\begin{columns}[T]
\column{0.48\textwidth}
De $x=a\sin\theta$:
\[
\sin\theta=\frac{x}{a}.
\]
El triángulo da
\[
\cos\theta=\frac{\sqrt{a^2-x^2}}{a}.
\]
Y por la rama principal,
\[
\theta=\arcsin\left(\frac{x}{a}\right).
\]
\column{0.48\textwidth}
\centering
\begin{tikzpicture}[scale=.82]
\coordinate (A) at (0,0);\coordinate (B) at (4,0);\coordinate (C) at (4,3);
\draw[thick] (A)--(B)--(C)--cycle;\draw (3.65,0)--(3.65,.35)--(4,.35);
\node at (2,-.35) {$\sqrt{a^2-x^2}$};\node[right] at (4,1.5) {$x$};\node[above left] at (2.05,1.65) {$a$};
\draw (.75,0) arc[start angle=0,end angle=36.87,radius=.75];\node at (.95,.32) {$\theta$};
\end{tikzpicture}
\end{columns}
\end{frame}

\begin{frame}{Ejemplo 1: regreso completo a $x$}
Partimos de
\[
I=\frac{a^2}{2}\theta+\frac{a^2}{2}\sin\theta\cos\theta+C.
\]
Sustituimos:
\begin{align*}
I
&=\frac{a^2}{2}\arcsin\left(\frac{x}{a}\right)
+\frac{a^2}{2}\left(\frac{x}{a}\right)
\left(\frac{\sqrt{a^2-x^2}}{a}\right)+C\\
&=\frac{a^2}{2}\arcsin\left(\frac{x}{a}\right)
+\frac{x}{2}\sqrt{a^2-x^2}+C.
\end{align*}
Por tanto,
\[
\boxed{\int\sqrt{a^2-x^2}\,dx=
\frac{x}{2}\sqrt{a^2-x^2}+
\frac{a^2}{2}\arcsin\left(\frac{x}{a}\right)+C.}
\]
\end{frame}

\begin{frame}{Ejemplo 1: ¿qué debemos aprender del procedimiento?}
No interesa únicamente la fórmula final. El mecanismo fue:
\[
\sqrt{a^2-x^2}
\xrightarrow{x=a\sin\theta}
a\cos\theta
\]
y simultáneamente
\[
dx\xrightarrow{x=a\sin\theta}a\cos\theta\,d\theta.
\]
Esto produjo
\[
a^2\int\cos^2\theta\,d\theta.
\]
Después el triángulo permitió reconstruir $\cos\theta$, mientras que la inversa del seno permitió reconstruir $\theta$.
\begin{alertblock}{Error frecuente}
Llegar a una respuesta en $\theta$ y detenerse. La integral original está en $x$; debemos regresar a $x$.
\end{alertblock}
\end{frame}

\section{Ejemplo 2: $a^2+x^2$}
\begin{frame}{Ejemplo 2: planteamiento}
Calcular
\[
\boxed{I=\int\frac{dx}{\sqrt{a^2+x^2}}},\qquad a>0.
\]
La estructura es
\[
a^2+x^2,
\]
que corresponde a
\[
1+\tan^2\theta=\sec^2\theta.
\]
Por ello,
\[
\boxed{x=a\tan\theta},\qquad
\boxed{dx=a\sec^2\theta\,d\theta}.
\]
\end{frame}

\begin{frame}{Ejemplo 2: transformar la raíz}
\begin{align*}
\sqrt{a^2+x^2}
&=\sqrt{a^2+a^2\tan^2\theta}\\
&=a\sqrt{1+\tan^2\theta}\\
&=a\sqrt{\sec^2\theta}\\
&=a|\sec\theta|.
\end{align*}
Tomamos
\[
-\frac{\pi}{2}<\theta<\frac{\pi}{2}.
\]
En este intervalo $\cos\theta>0$, por tanto $\sec\theta>0$ y
\[
\boxed{\sqrt{a^2+x^2}=a\sec\theta}.
\]
\end{frame}

\begin{frame}{Ejemplo 2: sustitución completa}
\begin{align*}
I
&=\int
\frac{\underbrace{a\sec^2\theta\,d\theta}_{dx}}
{\underbrace{a\sec\theta}_{\sqrt{a^2+x^2}}}\\
&=\int\sec\theta\,d\theta.
\end{align*}
Usamos la integral ya conocida:
\[
\int\sec\theta\,d\theta
=\ln|\sec\theta+\tan\theta|+C.
\]
Por tanto,
\[
\boxed{I=\ln|\sec\theta+\tan\theta|+C.}
\]
De nuevo, falta regresar a $x$.
\end{frame}

\begin{frame}{Ejemplo 2: triángulo y reconstrucción}
\begin{columns}[T]
\column{0.47\textwidth}
De
\[
\tan\theta=\frac{x}{a}
\]
tenemos opuesto $x$ y adyacente $a$. Entonces
\[
\text{hipotenusa}=\sqrt{a^2+x^2}
\]
y
\[
\sec\theta=\frac{\sqrt{a^2+x^2}}{a}.
\]
\column{0.50\textwidth}
\centering
\begin{tikzpicture}[scale=.86]
\coordinate (A) at (0,0);\coordinate (B) at (4,0);\coordinate (C) at (4,3);
\draw[thick] (A)--(B)--(C)--cycle;\draw (3.65,0)--(3.65,.35)--(4,.35);
\node at (2,-.35) {$a$};\node[right] at (4,1.5) {$x$};\node[above left] at (2.1,1.65) {$\sqrt{a^2+x^2}$};
\draw (.75,0) arc[start angle=0,end angle=36.87,radius=.75];\node at (.95,.32) {$\theta$};
\end{tikzpicture}
\end{columns}
\end{frame}

\begin{frame}{Ejemplo 2: regreso a $x$ y constante logarítmica}
\begin{align*}
I
&=\ln\left|
\frac{\sqrt{a^2+x^2}}{a}+\frac{x}{a}
\right|+C\\
&=\ln\left|\frac{x+\sqrt{a^2+x^2}}{a}\right|+C.
\end{align*}
Como $a>0$ es constante,
\[
\ln\left|\frac{x+\sqrt{a^2+x^2}}{a}\right|
=\ln|x+\sqrt{a^2+x^2}|-\ln a.
\]
Pero $-\ln a$ es una constante, así que se incorpora a $C$:
\[
\boxed{\int\frac{dx}{\sqrt{a^2+x^2}}
=\ln|x+\sqrt{a^2+x^2}|+C.}
\]
\end{frame}

\begin{frame}{Ejemplo 2: dos detalles que suelen causar dudas}
\begin{enumerate}
\item ¿Por qué aparece un logaritmo si empezamos con una raíz?
\[
\text{Porque la sustitución produjo }\int\sec\theta\,d\theta.
\]
\item ¿Por qué podemos ``desaparecer'' el factor $1/a$ dentro del logaritmo?
\[
\ln\left|\frac{F(x)}{a}\right|=\ln|F(x)|-\ln a,
\]
y $-\ln a$ es simplemente otra constante de integración.
\end{enumerate}
\begin{alertblock}{No confundir}
No estamos afirmando que $\ln|F/a|=\ln|F|$. Las primitivas difieren en una constante, y por eso representan la misma familia de antiderivadas.
\end{alertblock}
\end{frame}

\section{Ejemplo 3: $x^2-a^2$}
\begin{frame}{Ejemplo 3: planteamiento}
Calcular
\[
\boxed{I=\int\frac{dx}{\sqrt{x^2-a^2}}},\qquad a>0.
\]
La estructura es
\[
x^2-a^2,
\]
y la identidad correspondiente es
\[
\sec^2\theta-1=\tan^2\theta.
\]
Por ello,
\[
\boxed{x=a\sec\theta},\qquad
\boxed{dx=a\sec\theta\tan\theta\,d\theta}.
\]
\end{frame}

\begin{frame}{Ejemplo 3: transformar la raíz}
\begin{align*}
\sqrt{x^2-a^2}
&=\sqrt{a^2\sec^2\theta-a^2}\\
&=a\sqrt{\sec^2\theta-1}\\
&=a\sqrt{\tan^2\theta}\\
&=a|\tan\theta|.
\end{align*}
Para estudiar el caso $x\ge a$, elegimos
\[
0\le\theta<\frac{\pi}{2}.
\]
Entonces $\tan\theta\ge0$ y
\[
\boxed{\sqrt{x^2-a^2}=a\tan\theta}.
\]
\begin{observacion}
La integral real exige $|x|\ge a$. Aquí desarrollamos explícitamente la rama $x\ge a$; el valor absoluto en la forma final permite expresar la primitiva de manera apropiada en los intervalos de su dominio.
\end{observacion}
\end{frame}

\begin{frame}{Ejemplo 3: sustitución completa}
\begin{align*}
I
&=\int
\frac{\underbrace{a\sec\theta\tan\theta\,d\theta}_{dx}}
{\underbrace{a\tan\theta}_{\sqrt{x^2-a^2}}}\\
&=\int\sec\theta\,d\theta\\
&=\ln|\sec\theta+\tan\theta|+C.
\end{align*}
El patrón es parecido al ejemplo anterior, pero la reconstrucción geométrica es diferente porque ahora la relación inicial es
\[
\sec\theta=\frac{x}{a}.
\]
\end{frame}

\begin{frame}{Ejemplo 3: triángulo y reconstrucción}
\begin{columns}[T]
\column{0.47\textwidth}
Como
\[
\sec\theta=\frac{x}{a}
=\frac{\text{hipotenusa}}{\text{adyacente}},
\]
tomamos hipotenusa $x$ y adyacente $a$. Por Pitágoras,
\[
\text{opuesto}=\sqrt{x^2-a^2}.
\]
Así,
\[
\tan\theta=\frac{\sqrt{x^2-a^2}}{a}.
\]
\column{0.50\textwidth}
\centering
\begin{tikzpicture}[scale=.86]
\coordinate (A) at (0,0);\coordinate (B) at (4,0);\coordinate (C) at (4,3);
\draw[thick] (A)--(B)--(C)--cycle;\draw (3.65,0)--(3.65,.35)--(4,.35);
\node at (2,-.35) {$a$};\node[right] at (4,1.5) {$\sqrt{x^2-a^2}$};\node[above left] at (2.15,1.65) {$x$};
\draw (.75,0) arc[start angle=0,end angle=36.87,radius=.75];\node at (.95,.32) {$\theta$};
\end{tikzpicture}
\end{columns}
\end{frame}

\begin{frame}{Ejemplo 3: regreso a $x$}
\begin{align*}
I
&=\ln\left|
\frac{x}{a}+\frac{\sqrt{x^2-a^2}}{a}
\right|+C\\
&=\ln\left|
\frac{x+\sqrt{x^2-a^2}}{a}
\right|+C\\
&=\ln|x+\sqrt{x^2-a^2}|-\ln a+C.
\end{align*}
Absorbemos $-\ln a$ en la constante:
\[
\boxed{\int\frac{dx}{\sqrt{x^2-a^2}}
=\ln|x+\sqrt{x^2-a^2}|+C.}
\]
\end{frame}

\section{Comparación y práctica}
\begin{frame}{Comparación de los tres casos}
\centering
\renewcommand{\arraystretch}{1.5}
\begin{tabular}{>{\centering\arraybackslash}m{3.2cm}>{\centering\arraybackslash}m{3.1cm}>{\centering\arraybackslash}m{4.1cm}}
\toprule
Expresión & Sustitución & La raíz se transforma en\\
\midrule
$\sqrt{a^2-x^2}$ & $x=a\sin\theta$ & $a\cos\theta$\\
$\sqrt{a^2+x^2}$ & $x=a\tan\theta$ & $a\sec\theta$\\
$\sqrt{x^2-a^2}$ & $x=a\sec\theta$ & $a\tan\theta$\\
\bottomrule
\end{tabular}
\vspace{4mm}
La columna derecha no se memoriza por separado: se obtiene aplicando la identidad pitagórica correspondiente.
\end{frame}

\begin{frame}{Diagnóstico rápido: ¿qué sustitución elegiría?}
Sin integrar todavía, proponga una sustitución para cada caso:
\begin{enumerate}
\item $\displaystyle \int \frac{x^2}{\sqrt{25-x^2}}\,dx$
\item $\displaystyle \int \frac{dx}{(9+x^2)^{3/2}}$
\item $\displaystyle \int \sqrt{x^2-16}\,dx$
\item $\displaystyle \int \frac{dx}{x^2\sqrt{x^2-4}}$
\item $\displaystyle \int x^2\sqrt{36-x^2}\,dx$
\end{enumerate}
\pause
\[
\begin{array}{c|c}
1,5 & x=5\sin\theta\text{ o }x=6\sin\theta\\
2 & x=3\tan\theta\\
3 & x=4\sec\theta\\
4 & x=2\sec\theta
\end{array}
\]
\end{frame}

\begin{frame}{Errores frecuentes}
\begin{enumerate}
\item Elegir la sustitución por memoria sin verificar la identidad.
\item Cambiar la raíz pero olvidar transformar $dx$.
\item Escribir $\sqrt{u^2}=u$ sin considerar $|u|$.
\item No especificar un intervalo conveniente para $\theta$.
\item Construir mal el triángulo: confundir hipotenusa, opuesto y adyacente.
\item Terminar con una respuesta en $\theta$.
\item En expresiones logarítmicas, eliminar constantes sin justificar que se absorben en $C$.
\end{enumerate}
\end{frame}

\begin{frame}{Ejercicios para el pizarrón: nivel 1}
Determine primero la sustitución y dibuje el triángulo; después integre.
\begin{enumerate}
\item $\displaystyle \int\sqrt{9-x^2}\,dx$
\item $\displaystyle \int\frac{dx}{\sqrt{4+x^2}}$
\item $\displaystyle \int\frac{dx}{\sqrt{x^2-25}}$
\end{enumerate}
\vspace{2mm}
En cada ejercicio debe quedar explícito:
\[
\text{estructura}\to\text{identidad}\to\text{sustitución}\to dx\to\text{triángulo}.
\]
\end{frame}

\begin{frame}{Ejercicios para el pizarrón: nivel 2}
Ahora la expresión problemática está acompañada de otros factores:
\begin{enumerate}
\item $\displaystyle \int\frac{x^2}{\sqrt{16-x^2}}\,dx$
\item $\displaystyle \int\frac{x^2}{\sqrt{x^2+9}}\,dx$
\item $\displaystyle \int\frac{\sqrt{x^2-4}}{x}\,dx$
\item $\displaystyle \int\frac{dx}{(x^2+4)^{3/2}}$
\end{enumerate}
La sustitución trigonométrica simplifica la raíz, pero después todavía puede ser necesario utilizar identidades trigonométricas adicionales.
\end{frame}

\begin{frame}{Ejercicios para casa}
\begin{enumerate}
\item $\displaystyle \int\frac{dx}{\sqrt{49-x^2}}$
\item $\displaystyle \int\frac{x^2}{\sqrt{25-x^2}}\,dx$
\item $\displaystyle \int\sqrt{x^2+16}\,dx$
\item $\displaystyle \int\frac{dx}{(x^2+1)^{3/2}}$
\item $\displaystyle \int\sqrt{x^2-9}\,dx$
\item $\displaystyle \int\frac{dx}{x^2\sqrt{x^2-1}}$
\end{enumerate}
En los ejercicios 1--3, incluya explícitamente el triángulo utilizado para regresar a $x$.
\end{frame}

\begin{frame}{Mapa mental final}
\[
\boxed{
\begin{array}{c}
\text{¿Veo }a^2-x^2?\\
\Downarrow\\
x=a\sin\theta
\end{array}}
\qquad
\boxed{
\begin{array}{c}
\text{¿Veo }a^2+x^2?\\
\Downarrow\\
x=a\tan\theta
\end{array}}
\qquad
\boxed{
\begin{array}{c}
\text{¿Veo }x^2-a^2?\\
\Downarrow\\
x=a\sec\theta
\end{array}}
\]
\vspace{4mm}
Pero la pregunta más importante sigue siendo:
\[
\boxed{\text{¿qué identidad pitagórica hace funcionar esa sustitución?}}
\]
Si podemos responderla, ya no dependemos de memorizar una tabla.
\end{frame}

\begin{frame}{Conexión con las funciones trigonométricas inversas}
En el primer caso aparece naturalmente
\[
\theta=\arcsin\left(\frac{x}{a}\right),
\]
porque habíamos definido
\[
\sin\theta=\frac{x}{a}.
\]
En los otros casos, el triángulo permite reconstruir las funciones trigonométricas necesarias sin tener que escribir necesariamente $\theta$ mediante una inversa.
\vspace{2mm}
Esto enlaza la sustitución trigonométrica con las primitivas que contienen
\[
\arcsin x,\qquad \arctan x,\qquad \operatorname{arcsec}x,
\]
que estudiaremos a continuación desde el punto de vista de sus derivadas.
\end{frame}

\end{document}
